\section{The algorithm}

This section outlines the core concepts of the algorithm and its components, beginning with the successive approximation framework. Let $C = \max \{|c_{ij}|\colon (i,j) \in E\}$ denote the maximum possible cost in the considered network.

\subsection{The successive approximation framework}

The general procedure is outlined below.

Initially, the potentials $\pi$ and the flow $x$ are set to zero. To enforce $\epsilon$-optimality, we start with a sufficiently large $\epsilon$ that is at least as great as all arc costs in the network, i.e., $\epsilon = C$.

We then enter a loop. In each iteration, while $\epsilon > 1/n$, we halve $\epsilon$ and calculate the new $\epsilon$-optimal circulation $x$ and its potential $\pi$. Once $\epsilon < 1/n$, the optimal solution is found, according to \Cref{thm:optimal_epsi_circulation}.

The \textsc{Refine} operation is described later. To implement it, we will reuse several concepts established in the maximum flow problem.

\begin{algorithm}[H]
\caption{\textsc{SuccessiveApproximationAlgorithm}$(G,c, u)$}
\label{alg:succ_approx_framework}
    \begin{algorithmic}[1]
        \State $x \gets 0^m$ 
        \State $\pi \gets 0^n$
        \State $\epsilon \gets C$
        \While {$\epsilon > 1/n$}
        \State $\epsilon \gets \epsilon/2$
        \State $x, \pi \gets$ \Call{Refine}{$x,\pi,\epsilon$}
        \EndWhile
        \State\Return x
    \end{algorithmic}
\end{algorithm}


